\section*{Anexo 4.1-L --- Decisiones del numeral 16.1}
\addcontentsline{toc}{section}{Anexo 4.1-L --- Decisiones del numeral 16.1}
El detalle de decisiones del numeral 16.1 se presenta a continuación.

\begin{tablalafrox}{Decisiones de diseño del numeral 16.1}{tab:16-decisiones}{F{0.05} F{0.29} F{0.42} Y}{\cab{Nº} & \cab{Pregunta del caso} & \cab{Regla o componente propuesto} & \cab{Validación}}
1 & Entrega parcial e indicador de servicio & M6 distingue entregado completo, parcial y no entregado; M10 calcula OTIF por pedido completo y ventana acordada. & Criterio comercial \\
2 & Unidad de trazabilidad sanitaria & M9 traza GTIN y lote proveedor; M1/M5 vinculan cada movimiento a SSCC. & Unidad con Calidad \\
3 & Local cerrado y responsable del reintento & M6 registra causal y evidencia; la regla de reintento se parametriza y supervisa por Operaciones. & Plazo y autoridad \\
4 & Excursión térmica, decisión y bloqueo & M9 aplica umbral por producto; M5 bloquea localmente y Calidad autoriza liberar o descartar. & Umbrales con Calidad \\
5 & Identidad y sustitución de conductor externo & Portal Transportistas registra vínculo empresa--conductor--turno; OTP inicial y revocación al cambio; M6 conserva autor de cada POD. & Procedimiento transportista \\
6 & Efectivo y riesgo del dinero en ruta & M7 conserva cobro y rendición individual con causal de diferencia; Puelche define custodia y límite por turno. & Política de Tesorería \\
7 & Costo de servir y clientes no rentables & M10 calcula costo por entrega con ruta, tiempo, devoluciones y envases; Comercial decide medidas, no el algoritmo. & Criterio comercial \\
8 & Dos preventistas comprometen el mismo stock & M2 confirma reserva central; M3 deja pendiente el pedido sin conexión y comunica rechazo o sustitución al sincronizar. & Prioridad comercial \\
9 & Precio cambia antes del despacho & M3 conserva precio informado y versión de tarifa; una diferencia genera revalidación y aviso antes de confirmar. & Política de precios \\
10 & Control de envases retornables & M8 registra saldo por cliente, tipo y movimiento con evidencia; se concilia en la rendición. & Modalidad de cargo \\
11 & Maestro de productos ante cambios del proveedor & M1 ingresa equivalencia GTIN/formato; el maestro autorizado se sincroniza mediante ACL con ERP y preserva historial. & Dueño del maestro \\
12 & Devolución y DTE ya emitido & M8 registra devolución y causal; ACL solicita al ERP el documento tributario que corresponda, enlazado al original. & Regla tributaria \\
13 & Objeción sindical a cámaras y geolocalización & M12 usa telemetría de flota para ruta y costo; excluye cámaras en cabina y control de jornada por GPS. & Gestión laboral \\
14 & Destino del WMS 2013 & M1, M2 y M5 reemplazan sus capacidades en Etapa 1 por sitio y ola, con reversión controlada (ADR-08). & Corte por sitio \\
15 & Receptor distinto, imposibilidad o negativa a firmar & M6 identifica receptor y causal; conserva QR, fotografía o constancia de negativa según política, sin equiparar automáticamente POD con acuse tributario. & Validez de evidencia \\
16 & Conocimiento concentrado del planificador & M4 guarda restricciones, excepciones y decisiones del planificador; la transición incluye transferencia y prueba de rutas. & Aceptación Operaciones \\
\end{tablalafrox}

{\footnotesize Fuente: elaboración propia de LafroX a partir de Caso 02, numeral 16.1.\par}

% Protocolos del Anexo M y correspondencia del Anexo N.
